% =============================================================================
% Chapter 5 — The Virtual Graphics Controller ($A000–$A0BF)
% =============================================================================
\chapter{The Virtual Graphics Controller}

\epigraph{``Pixels are the atoms of the digital universe.  If you can control
them, you can build anything.''}{\textit{--- Overheard at a demo party, 1987}}

\section*{Introduction}

The Virtual Graphics Controller --- VGC for short --- is the heart of the
NovaVM's display system.  It manages three independent layers of visual output:
an 80$\times$50 text screen, a 320$\times$200 pixel bitmap with 16 colours, and
16 multicolour 16$\times$16 hardware sprites.  Current FPGA builds output a
720$\times$480 video signal with the 640$\times$400 Nova canvas centered inside
it.  All layers are composited at 60\,Hz into a single frame, with the
rendering order determined by the active display mode.

From the programmer's perspective, the VGC is a collection of memory-mapped
registers starting at \$A000.  The core registers at \$A000--\$A00F control
mode, colour, cursor, and status.  A command interface at \$A010--\$A01E
provides graphics primitives, sprite manipulation, copper (raster) effects, and
memory I/O.  Finally, 128 bytes of sprite registers at \$A040--\$A0BF give
direct access to each sprite's position, shape, and flags.

This chapter documents every register and command in the VGC.  It is the
largest chapter in the book for good reason: the VGC is where programs become
visible.


% ═══════════════════════════════════════════════════════════════════════════════
\section{Display Modes}
\label{sec:vgc-modes}
% ═══════════════════════════════════════════════════════════════════════════════

The VGC supports four display modes, selected by writing to \addr{A000}:

\begin{center}
\small
\begin{tabular}{@{}clp{22em}@{}}
\toprule
\textbf{Mode} & \textbf{Name} & \textbf{Description} \\
\midrule
0 & Text only     & 80$\times$50 text screen.  No graphics bitmap is drawn.
                     This is the default mode at power-on. \\
1 & Gfx over text & The graphics bitmap is drawn on top of the text layer.
                     Non-zero pixels in the bitmap obscure the text beneath. \\
2 & Text over gfx & The text layer is drawn on top of the graphics bitmap.
                     Characters with non-background colours obscure the bitmap. \\
3 & Gfx + sprites & Graphics bitmap and sprites only, no text layer.
                     Ideal for full-screen graphical games. \\
\bottomrule
\end{tabular}
\end{center}

\subsection*{Rendering Order}

Each frame at 60\,Hz proceeds in the following order:

\begin{enumerate}
  \item Background colour fill (the entire frame is cleared to \addr{A001}).
  \item Priority-0 sprites are drawn (behind everything).
  \item Text and/or graphics layers, in the order determined by the display mode.
  \item Priority-1 sprites are drawn (between text and graphics layers).
  \item Priority-2 sprites are drawn (in front of everything).
  \item Collision detection is performed and results are latched into
        \addr{A00B} and \addr{A00C}.
\end{enumerate}

\begin{lstlisting}
10 MODE 2 : REM text over graphics
20 GCLS   : REM clear bitmap
30 CIRCLE 160,100,80,1
40 PRINT "HELLO FROM TEXT LAYER"
\end{lstlisting}


% ═══════════════════════════════════════════════════════════════════════════════
\section{Core Registers (\$A000--\$A00F)}
\label{sec:vgc-core}
% ═══════════════════════════════════════════════════════════════════════════════

These sixteen registers control the fundamental display state: mode, colours,
cursor position, scrolling, and I/O.

% --- $A000 RegMode -----------------------------------------------------------

\mementry{A000}{RegMode}{R/W}{0}

Display mode selector.  Write a value from 0 to 3 to select the active display
mode (see Section~\ref{sec:vgc-modes}).

\begin{lstlisting}
POKE $A000,1 : REM graphics over text
\end{lstlisting}

\begin{notebox}
The BASIC command \texttt{MODE n} writes to this register.
\end{notebox}

% --- $A001 RegBgCol ----------------------------------------------------------

\mementry{A001}{RegBgCol}{R/W}{0}

Background colour (0--15).  This colour fills the entire screen behind all
layers.  It is also the ``paper'' colour for the text layer --- character cells
that are not covered by text are drawn in this colour.

\begin{lstlisting}
POKE $A001,6 : REM dark blue background
\end{lstlisting}

\begin{notebox}
\texttt{COLOR fg,bg} sets both \addr{A002} and \addr{A001} simultaneously.
\end{notebox}

% --- $A002 RegFgCol ----------------------------------------------------------

\mementry{A002}{RegFgCol}{R/W}{15}

Foreground/text colour (0--15).  Characters printed after changing this
register will appear in the new colour.

\begin{lstlisting}
POKE $A002,7 : REM white text
PRINT "NOW IN WHITE"
\end{lstlisting}

% --- $A003 RegCursorX --------------------------------------------------------

\mementry{A003}{RegCursorX}{R/W}{0}

Cursor column position (0--79).  Reading returns the current column; writing
moves the cursor.

\begin{lstlisting}
POKE $A003,40 : REM move cursor to column 40
PRINT "CENTERED-ISH"
\end{lstlisting}

\begin{notebox}
The BASIC command \texttt{LOCATE x,y} sets both \addr{A003} and \addr{A004}.
\end{notebox}

% --- $A004 RegCursorY --------------------------------------------------------

\mementry{A004}{RegCursorY}{R/W}{0}

Cursor row position (0--24).  Row 0 is the top of the screen.

\begin{lstlisting}
POKE $A004,12 : REM move cursor to row 12
\end{lstlisting}

% --- $A005 RegScrollX --------------------------------------------------------

\mementry{A005}{RegScrollX}{R/W}{0}

Horizontal scroll offset for the text layer.  Writing a non-zero value shifts
the text layer left by that many pixels.

\begin{lstlisting}
POKE $A005,4 : REM shift text 4 pixels left
\end{lstlisting}

% --- $A006 RegScrollY --------------------------------------------------------

\mementry{A006}{RegScrollY}{R/W}{0}

Vertical scroll offset for the text layer.  Shifts the text layer up by the
specified number of pixels.

\begin{lstlisting}
POKE $A006,2 : REM shift text 2 pixels up
\end{lstlisting}

% --- $A007 RegFont -----------------------------------------------------------

\mementry{A007}{RegFont}{R/W}{0}

Active font slot (0--7, 3 bits).  The VGC supports up to eight loaded font
definitions.  Slot 0 is the built-in system font.

\begin{lstlisting}
POKE $A007,1 : REM switch to font slot 1
\end{lstlisting}

% --- $A008 RegStatus ---------------------------------------------------------

\mementry{A008}{RegStatus}{R}{0}

Frame counter.  Increments by one at each vertical blank (60\,Hz) and wraps
from 255 back to 0.  This is the simplest way to synchronise with the display
refresh.

\begin{lstlisting}
10 F = PEEK($A008)
20 IF PEEK($A008) = F THEN 20
30 REM -- new frame has started --
\end{lstlisting}

\begin{tipbox}
The BASIC command \texttt{VSYNC} is equivalent to the wait loop above: it
blocks until the frame counter changes, then returns.
\end{tipbox}

% --- $A009 RegSpriteCount ----------------------------------------------------

\mementry{A009}{RegSpriteCount}{R}{0}

Returns the number of sprites currently enabled (0--16).

\begin{lstlisting}
PRINT PEEK($A009);" SPRITES ACTIVE"
\end{lstlisting}

% --- $A00A RegCursorEnable ---------------------------------------------------

\mementry{A00A}{RegCursorEnable}{R/W}{1}

Cursor visibility.  Any non-zero value displays the blinking cursor; zero
hides it.  Useful for games and full-screen applications where the cursor is
a distraction.

\begin{lstlisting}
POKE $A00A,0  : REM hide cursor
POKE $A00A,1  : REM show cursor
\end{lstlisting}

% --- $A00B RegColSt ----------------------------------------------------------

\mementry{A00B}{RegColSt}{R}{0}

Sprite-to-sprite collision bitmask.  After each frame, any sprite that
overlaps with another sprite has its corresponding bit set.  Bit $N$ represents
sprite $N$ (bit 0 = sprite 0, bit 7 = sprite 7).  For sprites 8--15 use the
BASIC \texttt{COLLISION()} function.

\begin{warningbox}
Reading this register \textbf{clears} it.  If you need to test multiple bits,
read the value into a variable first.
\end{warningbox}

\begin{lstlisting}
10 C = PEEK($A00B)
20 IF C AND 1 THEN PRINT "SPRITE 0 HIT SOMETHING"
30 IF C AND 2 THEN PRINT "SPRITE 1 HIT SOMETHING"
\end{lstlisting}

% --- $A00C RegColBg ----------------------------------------------------------

\mementry{A00C}{RegColBg}{R}{0}

Sprite-to-background collision bitmask.  Bit $N$ is set when sprite $N$
overlaps a non-zero pixel in the graphics bitmap.  Reading clears the
register, just like \addr{A00B}.

\begin{lstlisting}
10 B = PEEK($A00C)
20 IF B AND 1 THEN PRINT "SPRITE 0 HIT BACKGROUND"
\end{lstlisting}

% --- $A00D RegBorder ---------------------------------------------------------

\mementry{A00D}{RegBorder}{R/W}{11}

Border (frame) colour (0--15).  The border is the area outside the active
display area, if the host window is larger than the 320$\times$200 pixel
viewport.

\begin{lstlisting}
POKE $A00D,2 : REM red border
\end{lstlisting}

% --- $A00E RegCharOut --------------------------------------------------------

\mementry{A00E}{RegCharOut}{W}{---}

Character output port.  Writing a byte here places it on the screen at the
current cursor position and advances the cursor, exactly as if the character
had been PRINTed.  Control codes (carriage return, backspace, etc.)\ are
interpreted.

This is how BASIC's \texttt{PRINT} statement works internally --- every
character ultimately arrives here.

\begin{lstlisting}
10 POKE $A00E,72  : REM 'H'
20 POKE $A00E,73  : REM 'I'
30 POKE $A00E,13  : REM carriage return
\end{lstlisting}

% --- $A00F RegCharIn ---------------------------------------------------------

\mementry{A00F}{RegCharIn}{R}{---}

Character input port.  Returns the ASCII code of the next key in the keyboard
queue, or 0 if no key is pending.  This is how \texttt{INPUT} and \texttt{GET}
work internally --- the interpreter polls this register until a key arrives.

\begin{lstlisting}
10 K = PEEK($A00F)
20 IF K = 0 THEN 10
30 PRINT "YOU PRESSED: ";CHR$(K)
\end{lstlisting}


% ═══════════════════════════════════════════════════════════════════════════════
\section{Command Interface (\$A010--\$A01E)}
\label{sec:vgc-cmd}
% ═══════════════════════════════════════════════════════════════════════════════

The VGC command interface provides a rich set of graphics primitives, sprite
manipulation, copper (raster) effects, and memory I/O --- all triggered by
writing a command byte to the command register.

% --- $A010 RegCmd ------------------------------------------------------------

\mementry{A010}{RegCmd}{W}{---}

Command register.  Writing a command byte here executes the command
immediately.  Parameters must be loaded into P0--P14 (\$A011--\$A01F)
\emph{before} writing the command byte.

% --- $A011--$A01F RegP0--RegP14 ---------------------------------------------

\memrange{A011}{A01F}{RegP0--RegP14 (Parameter Registers)}{R/W}{0}

Fifteen parameter registers, addressed as \$A011 (P0) through \$A01F (P14).
Each command uses a different subset of these registers; unused parameters are
ignored.  Parameters are readable, so you can inspect the last-written values
at any time.

\begin{center}
\small
\begin{tabular}{@{}ccl@{}}
\toprule
\textbf{Address} & \textbf{Name} & \textbf{Offset from \$A011} \\
\midrule
\$A011 & P0  & +0 \\
\$A012 & P1  & +1 \\
\$A013 & P2  & +2 \\
\$A014 & P3  & +3 \\
\$A015 & P4  & +4 \\
\$A016 & P5  & +5 \\
\$A017 & P6  & +6 \\
\$A018 & P7  & +7 \\
\$A019 & P8  & +8 \\
\$A01A & P9  & +9 \\
\$A01B & P10 & +10 \\
\$A01C & P11 & +11 \\
\$A01D & P12 & +12 \\
\$A01E & P13 & +13 \\
\$A01F & P14 & +14 \\
\bottomrule
\end{tabular}
\end{center}

\begin{notebox}
Many parameters represent 16-bit values split across two consecutive registers
in little-endian order (low byte first, high byte second).  For example, an
X coordinate of 260 would be stored as P0=\$04 (low), P1=\$01 (high):
$260 = 256 \times 1 + 4$.
\end{notebox}


% ═══════════════════════════════════════════════════════════════════════════════
\subsection{Graphics Commands}
\label{sec:vgc-gfx-cmds}
% ═══════════════════════════════════════════════════════════════════════════════

% --- $01 CmdPlot -------------------------------------------------------------

\cmdentry{01}{CmdPlot}{Plot a single pixel}

Sets a pixel in the graphics bitmap.

\begin{center}
\small
\begin{tabular}{@{}lll@{}}
\toprule
\textbf{Param} & \textbf{Meaning} \\
\midrule
P0, P1 & X coordinate (16-bit LE, 0--319) \\
P2, P3 & Y coordinate (16-bit LE, 0--199) \\
P8     & Colour (0--15) \\
\bottomrule
\end{tabular}
\end{center}

\begin{lstlisting}
10 REM plot a red pixel at (100, 50)
20 POKE $A011,100 : POKE $A012,0 : REM X=100
30 POKE $A013,50  : POKE $A014,0 : REM Y=50
40 POKE $A019,2                   : REM colour=red
50 POKE $A010,$01                 : REM execute PLOT
\end{lstlisting}

\begin{notebox}
The BASIC command \texttt{PLOT x,y,c} does exactly this.
\end{notebox}

% --- $02 CmdUnplot -----------------------------------------------------------

\cmdentry{02}{CmdUnplot}{Erase a pixel (plot colour 0)}

Identical to CmdPlot but always uses colour 0 (the background colour),
effectively erasing the pixel.

\begin{center}
\small
\begin{tabular}{@{}ll@{}}
\toprule
\textbf{Param} & \textbf{Meaning} \\
\midrule
P0, P1 & X coordinate (16-bit LE) \\
P2, P3 & Y coordinate (16-bit LE) \\
\bottomrule
\end{tabular}
\end{center}

\begin{lstlisting}
10 POKE $A011,100 : POKE $A012,0
20 POKE $A013,50  : POKE $A014,0
30 POKE $A010,$02  : REM erase pixel at (100,50)
\end{lstlisting}

\begin{notebox}
BASIC equivalent: \texttt{UNPLOT x,y}.
\end{notebox}

% --- $03 CmdLine -------------------------------------------------------------

\cmdentry{03}{CmdLine}{Draw a line}

Draws a straight line between two endpoints using Bresenham's algorithm.

\begin{center}
\small
\begin{tabular}{@{}ll@{}}
\toprule
\textbf{Param} & \textbf{Meaning} \\
\midrule
P0, P1 & X1 (16-bit LE) \\
P2, P3 & Y1 (16-bit LE) \\
P4, P5 & X2 (16-bit LE) \\
P6, P7 & Y2 (16-bit LE) \\
P8     & Colour (0--15) \\
\bottomrule
\end{tabular}
\end{center}

\begin{lstlisting}
10 REM line from (0,0) to (319,199) in white
20 POKE $A011,0 : POKE $A012,0 : REM X1=0
30 POKE $A013,0 : POKE $A014,0 : REM Y1=0
40 POKE $A015,$3F : POKE $A016,$01 : REM X2=319
50 POKE $A017,$C7 : POKE $A018,$00 : REM Y2=199
60 POKE $A019,1 : REM colour=white
70 POKE $A010,$03 : REM execute LINE
\end{lstlisting}

\begin{notebox}
BASIC equivalent: \texttt{LINE x1,y1,x2,y2,c}.
\end{notebox}

% --- $04 CmdCircle -----------------------------------------------------------

\cmdentry{04}{CmdCircle}{Draw a circle outline}

Draws a circle using the midpoint circle algorithm.

\begin{center}
\small
\begin{tabular}{@{}ll@{}}
\toprule
\textbf{Param} & \textbf{Meaning} \\
\midrule
P0, P1 & Centre X (16-bit LE) \\
P2, P3 & Centre Y (16-bit LE) \\
P4, P5 & Radius (16-bit LE) \\
P8     & Colour (0--15) \\
\bottomrule
\end{tabular}
\end{center}

\begin{lstlisting}
10 REM circle at centre (160,100) radius 50, green
20 POKE $A011,160 : POKE $A012,0
30 POKE $A013,100 : POKE $A014,0
40 POKE $A015,50  : POKE $A016,0
50 POKE $A019,5   : REM colour=green
60 POKE $A010,$04  : REM execute CIRCLE
\end{lstlisting}

\begin{notebox}
BASIC equivalent: \texttt{CIRCLE x,y,r,c}.
\end{notebox}

% --- $05 CmdRect -------------------------------------------------------------

\cmdentry{05}{CmdRect}{Draw a rectangle outline}

Draws a rectangular outline from the top-left corner to the bottom-right
corner.

\begin{center}
\small
\begin{tabular}{@{}ll@{}}
\toprule
\textbf{Param} & \textbf{Meaning} \\
\midrule
P0, P1 & X1 top-left (16-bit LE) \\
P2, P3 & Y1 top-left (16-bit LE) \\
P4, P5 & X2 bottom-right (16-bit LE) \\
P6, P7 & Y2 bottom-right (16-bit LE) \\
P8     & Colour (0--15) \\
\bottomrule
\end{tabular}
\end{center}

\begin{lstlisting}
10 REM rectangle from (10,10) to (100,80) in cyan
20 POKE $A011,10 : POKE $A012,0
30 POKE $A013,10 : POKE $A014,0
40 POKE $A015,100 : POKE $A016,0
50 POKE $A017,80  : POKE $A018,0
60 POKE $A019,3   : REM colour=cyan
70 POKE $A010,$05  : REM execute RECT
\end{lstlisting}

\begin{notebox}
BASIC equivalent: \texttt{RECT x1,y1,x2,y2,c}.
\end{notebox}

% --- $06 CmdFill -------------------------------------------------------------

\cmdentry{06}{CmdFill}{Draw a filled rectangle}

Identical to CmdRect but fills the interior with the specified colour.

\begin{center}
\small
\begin{tabular}{@{}ll@{}}
\toprule
\textbf{Param} & \textbf{Meaning} \\
\midrule
P0, P1 & X1 top-left (16-bit LE) \\
P2, P3 & Y1 top-left (16-bit LE) \\
P4, P5 & X2 bottom-right (16-bit LE) \\
P6, P7 & Y2 bottom-right (16-bit LE) \\
P8     & Colour (0--15) \\
\bottomrule
\end{tabular}
\end{center}

\begin{lstlisting}
10 REM filled blue box
20 POKE $A011,0 : POKE $A012,0
30 POKE $A013,0 : POKE $A014,0
40 POKE $A015,$3F : POKE $A016,$01
50 POKE $A017,$C7 : POKE $A018,$00
60 POKE $A019,6 : REM colour=blue
70 POKE $A010,$06 : REM execute FILL
\end{lstlisting}

\begin{notebox}
BASIC equivalent: \texttt{FILL x1,y1,x2,y2,c}.
\end{notebox}

% --- $07 CmdGcls -------------------------------------------------------------

\cmdentry{07}{CmdGcls}{Clear graphics bitmap}

Clears the entire 320$\times$200 graphics bitmap to colour 0.  No parameters
are needed.

\begin{lstlisting}
POKE $A010,$07 : REM clear bitmap
\end{lstlisting}

\begin{notebox}
BASIC equivalent: \texttt{GCLS}.
\end{notebox}

% --- $08 CmdGcolor -----------------------------------------------------------

\cmdentry{08}{CmdGcolor}{Set default drawing colour}

Sets the default draw colour used by subsequent graphics commands when no
explicit colour is provided.

\begin{center}
\small
\begin{tabular}{@{}ll@{}}
\toprule
\textbf{Param} & \textbf{Meaning} \\
\midrule
P0 & Colour (0--15) \\
\bottomrule
\end{tabular}
\end{center}

\begin{lstlisting}
POKE $A011,5 : POKE $A010,$08 : REM default colour = green
\end{lstlisting}

\begin{notebox}
BASIC equivalent: \texttt{GCOLOR c}.
\end{notebox}

% --- $09 CmdPaint ------------------------------------------------------------

\cmdentry{09}{CmdPaint}{Flood fill from a point}

Performs a flood fill starting from the specified point, filling all
connected pixels of the same colour with the new colour.

\begin{center}
\small
\begin{tabular}{@{}ll@{}}
\toprule
\textbf{Param} & \textbf{Meaning} \\
\midrule
P0, P1 & X coordinate (16-bit LE) \\
P2, P3 & Y coordinate (16-bit LE) \\
P8     & Fill colour (0--15) \\
\bottomrule
\end{tabular}
\end{center}

\begin{lstlisting}
10 CIRCLE 160,100,50,1  : REM draw circle outline
20 REM flood fill the interior with yellow
30 POKE $A011,160 : POKE $A012,0
40 POKE $A013,100 : POKE $A014,0
50 POKE $A019,7   : REM colour=yellow
60 POKE $A010,$09  : REM execute PAINT
\end{lstlisting}

\begin{notebox}
BASIC equivalent: \texttt{PAINT x,y,c}.
\end{notebox}


% ═══════════════════════════════════════════════════════════════════════════════
\subsection{Sprite Commands}
\label{sec:vgc-spr-cmds}
% ═══════════════════════════════════════════════════════════════════════════════

Sprites are 16$\times$16 pixel objects with 4-bit colour per pixel (16 colours
including transparent).  The VGC supports 16 sprites numbered 0--15.  Sprite
shapes are stored in a separate shape RAM with 256 slots; each sprite register
block selects which shape slot to display.

% --- $10 CmdSprDef ----------------------------------------------------------

\cmdentry{10}{CmdSprDef}{Define a single sprite pixel}

Sets one pixel in a sprite's shape definition.

\begin{center}
\small
\begin{tabular}{@{}ll@{}}
\toprule
\textbf{Param} & \textbf{Meaning} \\
\midrule
P0 & Sprite number (0--15) \\
P1 & Pixel X within sprite (0--15) \\
P2 & Pixel Y within sprite (0--15) \\
P3 & Colour (0--15, 0 = transparent) \\
\bottomrule
\end{tabular}
\end{center}

\begin{lstlisting}
10 REM set pixel (7,7) of sprite 0 to red
20 POKE $A011,0   : REM sprite 0
30 POKE $A012,7   : REM x=7
40 POKE $A013,7   : REM y=7
50 POKE $A014,2   : REM colour=red
60 POKE $A010,$10  : REM execute SPRDEF
\end{lstlisting}

% --- $11 CmdSprRow -----------------------------------------------------------

\cmdentry{11}{CmdSprRow}{Load a sprite row (8 pixels)}

Loads eight pixels into one half of a sprite row.  Since sprites are 16 pixels
wide and each parameter byte holds one 4-bit colour value, two CmdSprRow
calls are needed per row --- but the BASIC \texttt{SPRITEDATA} command handles
this automatically.

\begin{center}
\small
\begin{tabular}{@{}ll@{}}
\toprule
\textbf{Param} & \textbf{Meaning} \\
\midrule
P0     & Sprite number (0--15) \\
P1     & Row index (0--15) \\
P2--P9 & Eight pixel colours (each 0--15) \\
\bottomrule
\end{tabular}
\end{center}

\begin{lstlisting}
10 REM load row 0 of sprite 0 with 8 red pixels
20 POKE $A011,0    : REM sprite 0
30 POKE $A012,0    : REM row 0
40 FOR I=0 TO 7
50   POKE $A013+I,2 : REM red
60 NEXT
70 POKE $A010,$11   : REM execute SPRROW
\end{lstlisting}

% --- $12 CmdSprClr ----------------------------------------------------------

\cmdentry{12}{CmdSprClr}{Clear a sprite's shape data}

Resets all pixels in the specified sprite's shape to transparent (colour 0).

\begin{center}
\small
\begin{tabular}{@{}ll@{}}
\toprule
\textbf{Param} & \textbf{Meaning} \\
\midrule
P0 & Sprite number (0--15) \\
\bottomrule
\end{tabular}
\end{center}

\begin{lstlisting}
POKE $A011,3 : POKE $A010,$12 : REM clear sprite 3
\end{lstlisting}

% --- $13 CmdSprCopy ---------------------------------------------------------

\cmdentry{13}{CmdSprCopy}{Copy a sprite's shape to another}

Duplicates the shape data from one sprite to another.  Useful for creating
animation frames or mirrored variants.

\begin{center}
\small
\begin{tabular}{@{}ll@{}}
\toprule
\textbf{Param} & \textbf{Meaning} \\
\midrule
P0 & Source sprite number \\
P1 & Destination sprite number \\
\bottomrule
\end{tabular}
\end{center}

\begin{lstlisting}
10 REM copy sprite 0's shape to sprite 1
20 POKE $A011,0 : POKE $A012,1
30 POKE $A010,$13
\end{lstlisting}

% --- $14 CmdSprPos -----------------------------------------------------------

\cmdentry{14}{CmdSprPos}{Set sprite position}

Positions a sprite at the given screen coordinates.  Sprite X is an unsigned
16-bit coordinate.  Sprite Y is an unsigned 8-bit coordinate over the
240-line sprite plane.

\begin{center}
\small
\begin{tabular}{@{}ll@{}}
\toprule
\textbf{Param} & \textbf{Meaning} \\
\midrule
P0       & Sprite number (0--15) \\
P1, P2   & X position (unsigned 16-bit LE) \\
P3       & Y position (unsigned 8-bit) \\
P4       & Reserved; write 0 \\
\bottomrule
\end{tabular}
\end{center}

\begin{lstlisting}
10 REM position sprite 2 at (160, 100)
20 POKE $A011,2    : REM sprite 2
30 POKE $A012,160  : POKE $A013,0 : REM X=160
40 POKE $A014,100  : POKE $A015,0 : REM Y=100, P4 reserved
50 POKE $A010,$14   : REM execute SPRPOS
\end{lstlisting}

\begin{notebox}
The BASIC command \texttt{SPRITE n,x,y} sets position and enables the sprite
in one step.
\end{notebox}

% --- $15 CmdSprEna -----------------------------------------------------------

\cmdentry{15}{CmdSprEna}{Enable (show) a sprite}

Makes the sprite visible on screen.  The sprite must have shape data and a
position set, or it will appear as a transparent rectangle at (0,0).

\begin{center}
\small
\begin{tabular}{@{}ll@{}}
\toprule
\textbf{Param} & \textbf{Meaning} \\
\midrule
P0 & Sprite number (0--15) \\
\bottomrule
\end{tabular}
\end{center}

\begin{lstlisting}
POKE $A011,0 : POKE $A010,$15 : REM enable sprite 0
\end{lstlisting}

% --- $16 CmdSprDis ----------------------------------------------------------

\cmdentry{16}{CmdSprDis}{Disable (hide) a sprite}

Removes the sprite from the display.  Its shape and position data are
preserved.

\begin{center}
\small
\begin{tabular}{@{}ll@{}}
\toprule
\textbf{Param} & \textbf{Meaning} \\
\midrule
P0 & Sprite number (0--15) \\
\bottomrule
\end{tabular}
\end{center}

\begin{lstlisting}
POKE $A011,0 : POKE $A010,$16 : REM disable sprite 0
\end{lstlisting}

\begin{notebox}
BASIC equivalent: \texttt{SPRITE n,OFF}.
\end{notebox}

% --- $17 CmdSprFlip ----------------------------------------------------------

\cmdentry{17}{CmdSprFlip}{Set sprite flip flags}

Mirrors the sprite horizontally and/or vertically.

\begin{center}
\small
\begin{tabular}{@{}ll@{}}
\toprule
\textbf{Param} & \textbf{Meaning} \\
\midrule
P0 & Sprite number (0--15) \\
P1 & Flip flags: bit 0 = X flip, bit 1 = Y flip \\
\bottomrule
\end{tabular}
\end{center}

\begin{lstlisting}
10 POKE $A011,0   : REM sprite 0
20 POKE $A012,1   : REM flip horizontally
30 POKE $A010,$17
\end{lstlisting}

Values for P1:
\begin{itemize}
  \item 0 = no flip
  \item 1 = horizontal flip (mirror left-right)
  \item 2 = vertical flip (mirror top-bottom)
  \item 3 = both flips (180-degree rotation)
\end{itemize}

% --- $18 CmdSprPri ----------------------------------------------------------

\cmdentry{18}{CmdSprPri}{Set sprite priority}

Controls the sprite's draw order relative to the text and graphics layers.

\begin{center}
\small
\begin{tabular}{@{}ll@{}}
\toprule
\textbf{Param} & \textbf{Meaning} \\
\midrule
P0 & Sprite number (0--15) \\
P1 & Priority (0--2) \\
\bottomrule
\end{tabular}
\end{center}

Priority values:
\begin{itemize}
  \item 0 = behind all layers (drawn first)
  \item 1 = between text and graphics layers
  \item 2 = in front of all layers (drawn last)
\end{itemize}

\begin{lstlisting}
10 POKE $A011,0   : REM sprite 0
20 POKE $A012,2   : REM in front of everything
30 POKE $A010,$18
\end{lstlisting}


% ═══════════════════════════════════════════════════════════════════════════════
\subsection{Memory I/O Commands}
\label{sec:vgc-mem-cmds}
% ═══════════════════════════════════════════════════════════════════════════════

These commands allow direct read/write access to the VGC's internal memory
spaces, which are not mapped into the CPU's address space.

% --- $19 CmdMemRead ----------------------------------------------------------

\cmdentry{19}{CmdMemRead}{Read from VGC memory}

Reads a byte from one of the VGC's internal memory spaces and returns it in
P3.

\begin{center}
\small
\begin{tabular}{@{}ll@{}}
\toprule
\textbf{Param} & \textbf{Meaning} \\
\midrule
P0     & Memory space selector \\
P1, P2 & Address within space (16-bit LE) \\
P3     & \textit{(output)} Data byte read \\
P4     & Flags: bit 0 = auto-increment address \\
\bottomrule
\end{tabular}
\end{center}

Memory space values:
\begin{itemize}
  \item \$01 = Character RAM (4000 bytes, character codes)
  \item \$02 = Colour RAM (4000 bytes, per-cell colour attributes)
  \item \$03 = Graphics bitmap (64000 bytes, 320$\times$200 pixels)
  \item \$04 = Sprite shape RAM (256 slots $\times$ 128 bytes = 32768 bytes)
  \item \$06 = Tile data RAM (32768 bytes)
\end{itemize}

\begin{lstlisting}
10 REM read the character at screen position 0
20 POKE $A011,1  : REM space = character RAM
30 POKE $A012,0  : POKE $A013,0 : REM address = 0
40 POKE $A015,0  : REM no auto-increment
50 POKE $A010,$19 : REM execute MEMREAD
60 PRINT "CHAR CODE: ";PEEK($A014)
\end{lstlisting}

\begin{tipbox}
Set P4 bit 0 to enable auto-increment mode.  After each read, the address in
P1/P2 is automatically advanced by one, making it efficient to read
consecutive bytes in a loop without rewriting the address registers.
\end{tipbox}

% --- $1A CmdMemWrite ---------------------------------------------------------

\cmdentry{1A}{CmdMemWrite}{Write to VGC memory}

Writes a byte to one of the VGC's internal memory spaces.

\begin{center}
\small
\begin{tabular}{@{}ll@{}}
\toprule
\textbf{Param} & \textbf{Meaning} \\
\midrule
P0     & Memory space selector (same as CmdMemRead) \\
P1, P2 & Address within space (16-bit LE) \\
P3     & Data byte to write \\
P4     & Flags: bit 0 = auto-increment address \\
\bottomrule
\end{tabular}
\end{center}

\begin{lstlisting}
10 REM write 'A' (65) to screen position 0
20 POKE $A011,1   : REM space = character RAM
30 POKE $A012,0   : POKE $A013,0 : REM address = 0
40 POKE $A014,65  : REM data = 'A'
50 POKE $A015,0   : REM no auto-increment
60 POKE $A010,$1A  : REM execute MEMWRITE
\end{lstlisting}


% ═══════════════════════════════════════════════════════════════════════════════
\subsection{Copper (Raster) Commands}
\label{sec:vgc-copper-cmds}
% ═══════════════════════════════════════════════════════════════════════════════

The copper is a raster-synchronised coprocessor inspired by the Amiga's
Copper.  It executes a list of register-write or IRQ events keyed to
specific raster positions, allowing per-scanline colour changes, split-screen
effects, raster bars, and beam-position CPU notifications --- all without
CPU polling.

FPGA builds support 8 copper lists, each holding up to 32 events.  Simulation
builds may expose more lists for software compatibility.  One list is
\emph{active} (being executed each frame) and one is the \emph{target} (being
edited).  This double-buffering allows glitch-free list updates.

% --- $1B CmdCopperAdd -------------------------------------------------------

\cmdentry{1B}{CmdCopperAdd}{Add a copper event}

Adds an event to the target copper list.

\begin{center}
\small
\begin{tabular}{@{}ll@{}}
\toprule
\textbf{Param} & \textbf{Meaning} \\
\midrule
P0, P1 & X position (16-bit LE, 0--319) \\
P2     & Y position (0--199) \\
P3, P4 & Target register or pseudo-register \\
P5     & Value to write to the register \\
\bottomrule
\end{tabular}
\end{center}

\begin{lstlisting}
10 REM change background colour at scanline 100
20 POKE $A011,0   : POKE $A012,0  : REM x = 0
30 POKE $A013,100                  : REM y = 100
40 POKE $A014,$01 : POKE $A015,$A0 : REM target = $A001 (BgCol)
50 POKE $A016,2                    : REM value = red
60 POKE $A010,$1B                  : REM add event
\end{lstlisting}

\begin{notebox}
The BASIC command \texttt{COPPER} provides a higher-level interface to the
copper system.
\end{notebox}

\begin{tipbox}
Target pseudo-register \texttt{\$FE} raises VGC IRQ pending bits from P5
instead of writing a display register.  This lets the copper request CPU
attention at arbitrary beam positions without a fixed C64-style raster IRQ
compare register.
\end{tipbox}

% --- $1C CmdCopperClear ------------------------------------------------------

\cmdentry{1C}{CmdCopperClear}{Clear the target copper list}

Removes all events from the currently selected target copper list.  No
parameters.

\begin{lstlisting}
POKE $A010,$1C : REM clear target list
\end{lstlisting}

% --- $1D CmdCopperEnable -----------------------------------------------------

\cmdentry{1D}{CmdCopperEnable}{Enable copper execution}

Starts executing the active copper list each frame.  No parameters.

\begin{lstlisting}
POKE $A010,$1D : REM enable copper
\end{lstlisting}

% --- $1E CmdCopperDisable ----------------------------------------------------

\cmdentry{1E}{CmdCopperDisable}{Disable copper execution}

Stops copper execution.  The list contents are preserved.  No parameters.

\begin{lstlisting}
POKE $A010,$1E : REM disable copper
\end{lstlisting}

% --- $1F CmdSysReset ---------------------------------------------------------

\cmdentry{1F}{CmdSysReset}{Full system reset}

Performs a complete visual and audio reset: clears the screen, silences all
audio, disables all sprites, resets colours to defaults.  No parameters.

\begin{lstlisting}
POKE $A010,$1F : REM full reset
\end{lstlisting}

\begin{notebox}
BASIC equivalent: \texttt{RESET}.
\end{notebox}

% --- $20 CmdCopperList -------------------------------------------------------

\cmdentry{20}{CmdCopperList}{Select copper list for editing}

Selects which copper list is the \emph{target} for subsequent CmdCopperAdd
and CmdCopperClear operations.

\begin{center}
\small
\begin{tabular}{@{}ll@{}}
\toprule
\textbf{Param} & \textbf{Meaning} \\
\midrule
P0 & List index (0--127) \\
\bottomrule
\end{tabular}
\end{center}

\begin{lstlisting}
POKE $A011,1 : POKE $A010,$20 : REM edit list 1
\end{lstlisting}

% --- $21 CmdCopperUse --------------------------------------------------------

\cmdentry{21}{CmdCopperUse}{Activate a copper list}

Switches the active copper list to the specified index.  The switch takes
effect at the next vertical blank.

\begin{center}
\small
\begin{tabular}{@{}ll@{}}
\toprule
\textbf{Param} & \textbf{Meaning} \\
\midrule
P0 & List index (0--127) \\
\bottomrule
\end{tabular}
\end{center}

\begin{lstlisting}
POKE $A011,1 : POKE $A010,$21 : REM activate list 1
\end{lstlisting}

% --- $22 CmdCopperListEnd ----------------------------------------------------

\cmdentry{22}{CmdCopperListEnd}{Reset target to active list}

Sets the target copper list back to whichever list is currently active.
Useful after finishing edits on a different list.  No parameters.

\begin{lstlisting}
POKE $A010,$22 : REM target = active list
\end{lstlisting}


% ═══════════════════════════════════════════════════════════════════════════════
\section{VGC IRQ Block (\$A0F0--\$A0FF)}
\label{sec:vgc-irq}
% ═══════════════════════════════════════════════════════════════════════════════

\mementry{A0F0}{RegIrqEnable}{R/W}{0}

Enabled IRQ source mask.  Bit 0 is vblank.  Bits 1--4 are copper/user event
sources raised by copper pseudo-register \texttt{\$FE}.

\mementry{A0F1}{RegIrqStatus}{R/W1C}{0}

Pending IRQ source mask.  The CPU IRQ line is asserted while
\texttt{RegIrqStatus \& RegIrqEnable} is non-zero.  Write one bits to clear
handled sources.

\mementry{A0F2}{RegIrqForce}{W}{---}

Sets enabled pending bits for diagnostics and test programs.

\mementry{A0F3}{RegIrqValid}{R}{\$1F}

Returns the implemented IRQ source mask.

\begin{lstlisting}
POKE $A0F0,1 : REM enable vblank IRQ source
POKE $A0F1,1 : REM acknowledge vblank IRQ source
\end{lstlisting}

\begin{tipbox}
Nova does not provide a fixed C64-style raster-line compare register.  Use
the copper for raster-position effects, and use copper IRQ events only when
the CPU really needs to be notified.
\end{tipbox}


% ═══════════════════════════════════════════════════════════════════════════════
\section{Help System Registers (\$A020--\$A030)}
\label{sec:vgc-help}
% ═══════════════════════════════════════════════════════════════════════════════

The VGC includes a built-in help system accessible through a small register
block.

\mementry{A020}{RegHelp}{W}{---}

Help command register.  Write one of the following values:
\begin{itemize}
  \item \$01 --- Display the help index (list of available topics).
  \item \$02 --- Search for a keyword.  The search query must be written to
        the search buffer at \addr{A021} first.
\end{itemize}

\begin{lstlisting}
POKE $A020,1 : REM show help index
\end{lstlisting}

\memrange{A021}{A030}{HelpSearchBuffer}{R/W}{0}

A 16-byte buffer for the help search query.  Write a null-terminated ASCII
string here, then write \$02 to \addr{A020} to execute the search.

\begin{lstlisting}
10 REM search for "PLOT" in help
20 POKE $A021,ASC("P")
30 POKE $A022,ASC("L")
40 POKE $A023,ASC("O")
50 POKE $A024,ASC("T")
60 POKE $A025,0         : REM null terminator
70 POKE $A020,2         : REM search
\end{lstlisting}


% ═══════════════════════════════════════════════════════════════════════════════
\section{Compiler Controller (\$A031--\$A03E)}
\label{sec:vgc-compiler}
% ═══════════════════════════════════════════════════════════════════════════════

The NovaVM includes an on-board NCC compiler.  These registers provide
access to compilation commands, status, and error reporting.

\mementry{A031}{CmpCmd}{W}{---}

Compiler command register.  Write one of:
\begin{itemize}
  \item \$01 --- Compile source code from XRAM at the address in \addr{A03B}--\addr{A03D}.
  \item \$02 --- Populate error fields for the error at index \addr{A03A}.
  \item \$03 --- Populate warning fields for the warning at index \addr{A03A}.
\end{itemize}

\mementry{A032}{CmpStatus}{R}{0}

Compiler status:
\begin{itemize}
  \item \$00 = idle
  \item \$01 = compiling
  \item \$02 = compilation successful
  \item \$03 = compilation failed with errors
\end{itemize}

\mementry{A033}{CmpErrCount}{R}{0}

Number of errors from the last compilation.

\mementry{A034}{CmpWarnCount}{R}{0}

Number of warnings from the last compilation.

\memrange{A035}{A036}{CmpErrLine}{R}{0}

16-bit line number (little-endian) of the current error.

\memrange{A037}{A038}{CmpCodeSize}{R}{0}

16-bit size (little-endian) of the compiled code in bytes.

\mementry{A039}{CmpErrMsg}{R}{---}

Read this register repeatedly to retrieve the error message one byte at a
time.  Returns 0 when the message is exhausted.

\mementry{A03A}{CmpErrIdx}{W}{---}

Write an error or warning index here (0-based), then issue command \$02 or
\$03 to \addr{A031} to populate the error fields.

\memrange{A03B}{A03D}{CmpSrcAddr}{W}{---}

24-bit little-endian XRAM address of the source code to compile.  Set this
before issuing the compile command.


% ═══════════════════════════════════════════════════════════════════════════════
\section{ROM Swap Register (\$A03F)}
\label{sec:vgc-romswap}
% ═══════════════════════════════════════════════════════════════════════════════

\mementry{A03F}{RegRomSwap}{W}{---}

Controls which ROM image is active in the \$C000--\$FFFF region.

\begin{itemize}
  \item \$01 --- Switch to NCC ROM (native code compiler).
  \item \$02 --- Switch to EhBASIC ROM (the default).
  \item \$03 --- Activate the NCC editor.
\end{itemize}

\begin{warningbox}
Switching ROMs while BASIC is running will crash the system.  This register is
intended for use by the boot loader and compiler infrastructure, not by user
programs.
\end{warningbox}


% ═══════════════════════════════════════════════════════════════════════════════
\section{Sprite Registers (\$A040--\$A0BF)}
\label{sec:vgc-sprite-regs}
% ═══════════════════════════════════════════════════════════════════════════════

The 128 bytes from \$A040 to \$A0BF are divided into 16 blocks of 8 bytes
each, one block per sprite.  These registers provide direct read/write access
to each sprite's position, shape slot, flags, priority, and transparent
colour --- bypassing the command interface entirely.

The base address of sprite $N$ is:

\begin{center}
\texttt{Base = \$A040 + $N$ $\times$ 8}
\end{center}

\begin{center}
\small
\begin{tabular}{@{}clll@{}}
\toprule
\textbf{Offset} & \textbf{Name} & \textbf{R/W} & \textbf{Description} \\
\midrule
+0 & SprRegXLo       & R/W & X position low byte \\
+1 & SprRegXHi       & R/W & X position high byte \\
+2 & SprRegYLo       & R/W & Y position (unsigned byte) \\
+3 & SprRegYHi       & R/W & Reserved; reads as 0 \\
+4 & SprRegShape     & R/W & Shape slot index (0--255) \\
+5 & SprRegFlags     & R/W & Flags (see below) \\
+6 & SprRegPriority  & R/W & Draw priority (0--2) \\
+7 & SprRegTransColor & R/W & Transparent colour (default 0) \\
\bottomrule
\end{tabular}
\end{center}

\subsection*{SprRegFlags (offset +5)}

\bitfield{ENA}{---}{---}{---}{---}{---}{YFLP}{XFLP}

\begin{itemize}
  \item \textbf{Bit 7 (ENA)}: Enable.  Set to show the sprite, clear to hide it.
  \item \textbf{Bit 1 (YFLP)}: Vertical flip.
  \item \textbf{Bit 0 (XFLP)}: Horizontal flip.
\end{itemize}

Bits 2--6 are unused and should be written as zero.

\subsection*{Sprite Register Addresses}

For quick reference, here are the base addresses for all 16 sprites:

\begin{center}
\small
\begin{tabular}{@{}cccc@{}}
\toprule
\textbf{Sprite} & \textbf{Base} & \textbf{Sprite} & \textbf{Base} \\
\midrule
0  & \$A040 & 8  & \$A080 \\
1  & \$A048 & 9  & \$A088 \\
2  & \$A050 & 10 & \$A090 \\
3  & \$A058 & 11 & \$A098 \\
4  & \$A060 & 12 & \$A0A0 \\
5  & \$A068 & 13 & \$A0A8 \\
6  & \$A070 & 14 & \$A0B0 \\
7  & \$A078 & 15 & \$A0B8 \\
\bottomrule
\end{tabular}
\end{center}

\subsection*{Sprite Register Access}

NovaVM exposes sprite registers for low-level code, but NovaBASIC programs
should use the sprite commands and functions instead.  The FPGA buffers sprite
position, shape, flags, priority, and transparency changes and publishes them at
the frame boundary, so immediate writes do not tear visible sprites.

\begin{lstlisting}
10 REM position sprite 3 at (100, 50)
20 SPRITE 3,ON
30 SPRITE 3,100,50
\end{lstlisting}

\begin{lstlisting}
10 REM read sprite 5's current position
20 X = SPRITEX(5)
30 Y = SPRITEY(5)
40 PRINT "SPRITE 5 AT ";X;",";Y
\end{lstlisting}

\begin{notebox}
Y high is reserved.  Sprite X is an unsigned 16-bit value; sprite Y is an
unsigned 8-bit value over the 0--239 sprite plane.
\end{notebox}

\begin{lstlisting}
10 REM enable sprite 0 with horizontal flip
20 POKE $A045, $80 OR $01 : REM ENA + XFLP
\end{lstlisting}

\begin{lstlisting}
10 REM set sprite 0 to draw behind everything
20 POKE $A046, 0 : REM priority = 0
\end{lstlisting}

\begin{notebox}
The transparent colour register (offset +7) defaults to 0, meaning black
pixels in the sprite shape are treated as transparent.  Change this if your
sprite needs to contain black pixels --- for example, set it to 15 to make
light grey transparent instead.
\end{notebox}


% ═══════════════════════════════════════════════════════════════════════════════
\section*{Summary}
% ═══════════════════════════════════════════════════════════════════════════════

The VGC occupies addresses \$A000--\$A0BF and provides the complete visual
output system for the NovaVM.  The table below summarises the register map:

\begin{center}
\small
\begin{tabular}{@{}lll@{}}
\toprule
\textbf{Range} & \textbf{Usage} & \textbf{Access} \\
\midrule
\$A000--\$A00F  & Core registers (mode, colour, cursor, I/O) & R/W \\
\$A010          & Command register                           & W \\
\$A011--\$A01F  & Parameter registers P0--P14                & R/W \\
\$A020--\$A030  & Help system                                & R/W \\
\$A031--\$A03E  & Compiler controller                        & R/W \\
\$A03F          & ROM swap                                   & W \\
\$A040--\$A0BF  & Sprite registers (16 $\times$ 8 bytes)     & R/W \\
\$A0F0--\$A0FF  & VGC IRQ block                              & R/W \\
\bottomrule
\end{tabular}
\end{center}

With 9 graphics commands, 9 sprite commands, 2 memory I/O commands, 8 copper
commands, and 128 bytes of direct sprite registers, the VGC packs a remarkable
amount of capability into a small address range.  Master this chapter, and
you have mastered the display.
